% unit: u11
% sources: sources/099-the-four-fundamental-subspaces.txt, sources/157-the-four-fundamental-subspaces.txt, sources/158-mit18-06scf11-ses1-10sum-pdf-download-file.txt, sources/159-computing-the-four-fundamental-subspaces.txt, sources/160-mit18-06scf11-ses1-10prob-pdf-download-file.txt, sources/161-mit18-06scf11-ses1-10sol-pdf-download-file.txt
\section{The four fundamental subspaces}\label{sec:u11}

So far each matrix has come with two subspaces, the column space and the nullspace. Transposing
the matrix gives two more, and the four together are the heart of this approach to linear algebra.
This section finds a basis and the dimension of each, all from one elimination $EA=R$, and proves the
first great count: the row space and the column space have the same dimension $r$. That fact
already corrects the slip of Section~\ref{sec:u10}, where $(3,3,8)$ was proposed as independent of
$(1,1,2)$ and $(2,2,5)$: the matrix with those columns has two equal rows, so its row space is only
two-dimensional, and therefore so is its column space.

\subsection{Two subspaces in \texorpdfstring{\(\R^n\)}{Rn}, two in \texorpdfstring{\(\R^m\)}{Rm}}

\begin{definition}[The four fundamental subspaces]\label{def:u11-four}
Let $A$ be an $m\times n$ matrix.
\begin{itemize}
\item The \emph{column space} $\Col{A}\subseteq\R^m$: all combinations of the columns.
\item The \emph{nullspace} $\Nul{A}\subseteq\R^n$: all solutions of $A\vb{x}=\vb{0}$.
\item The \emph{row space} $\Col{A\T}\subseteq\R^n$: all combinations of the rows of $A$, written as
column vectors, i.e.\ the column space of $A\T$.
\item The \emph{left nullspace} $\Nul{A\T}\subseteq\R^m$: all solutions of $A\T\vb{y}=\vb{0}$.
\end{itemize}
\end{definition}

Transposing $A\T\vb{y}=\vb{0}$ gives $\vb{y}\T A=\vb{0}\T$: a row vector multiplying $A$ from the
\emph{left} gives the zero row. That explains the name. A vector $\vb{y}$ in the left nullspace lists
the coefficients of a combination of the rows of $A$ that produces the zero row, just as a vector in
$\Nul{A}$ lists a combination of the columns that produces the zero column. These are exactly the
vectors that appeared in the solvability test of Theorem~\ref{thm:u09-solvable}.
Figure~\ref{fig:u11-big} shows where the four spaces live.

\begin{figure}[ht]\centering
\begin{tikzpicture}[>=Stealth,font=\small]
  \begin{scope}
    \coordinate (On) at (0,0);
    \filldraw[fill=blue!8,draw=blue!60!black] (On) -- ++(45:1.8) -- ++(135:1.8) -- ++(225:1.8) -- cycle;
    \node at (0,1.27) {dim $r$};
    \node[anchor=east,align=right] at (-1.35,1.6) {row space\\$\Col{A\T}$};
    \filldraw[fill=gray!12,draw=gray!50!black] (On) -- ++(-45:1.55) -- ++(-135:1.55) -- ++(135:1.55) -- cycle;
    \node[font=\footnotesize] at (0,-1.1) {dim $n-r$};
    \node[anchor=east,align=right] at (-1.2,-1.4) {nullspace\\$\Nul{A}$};
    \fill (On) circle (1.5pt);
    \node at (0,-2.6) {$\R^n$};
  \end{scope}
  \begin{scope}[xshift=7.2cm]
    \coordinate (Om) at (0,0);
    \filldraw[fill=orange!12,draw=orange!70!black] (Om) -- ++(45:1.8) -- ++(135:1.8) -- ++(225:1.8) -- cycle;
    \node at (0,1.27) {dim $r$};
    \node[anchor=west,align=left] at (1.35,1.6) {column space\\$\Col{A}$};
    \filldraw[fill=gray!12,draw=gray!50!black] (Om) -- ++(-45:1.35) -- ++(-135:1.35) -- ++(135:1.35) -- cycle;
    \node[font=\footnotesize] at (0,-0.95) {dim $m-r$};
    \node[anchor=west,align=left] at (1.05,-1.25) {left nullspace\\$\Nul{A\T}$};
    \fill (Om) circle (1.5pt) node[right=2pt] {$\vb{0}$};
    \node at (0,-2.6) {$\R^m$};
  \end{scope}
  \draw[->,thick] (0.75,1.75) to[bend left=12] node[above] {$\vb{x}\mapsto A\vb{x}$} (6.45,1.75);
  \draw[->,thick] (1.15,-1.1) to[bend left=6] node[below,pos=0.45] {$A\vb{x}=\vb{0}$} (7.1,-0.05);
\end{tikzpicture}
\caption{The four fundamental subspaces of an $m\times n$ matrix of rank $r$. The row space and
nullspace live in $\R^n$ and their dimensions add to $n$; the column space and left nullspace live
in $\R^m$ and their dimensions add to $m$. $A$ sends every input into the column space and sends the
nullspace to $\vb{0}$.}
\label{fig:u11-big}
\end{figure}

For the column space and the nullspace the work is already done (Theorems~\ref{thm:u10-col}
and~\ref{thm:u10-null}): the pivot columns of $A$ are a basis of $\Col{A}$, of dimension $r$, and
the $n-r$ special solutions are a basis of $\Nul{A}$. One could treat $A\T$ the same way, but that
means a second elimination. The point is that the elimination already done on $A$ suffices.

\subsection{The row space}

\begin{theorem}[Row space]\label{thm:u11-row}
If $E$ is invertible and $EA=R$, then $A$ and $R$ have the same row space. In particular the $r$
nonzero rows of $R=\operatorname{rref}(A)$ form a basis of $\Col{A\T}$, and $\dim\Col{A\T}=r$.
\end{theorem}
\begin{proof}
Row $i$ of $R=EA$ is $\sum_k e_{ik}(\text{row }k\text{ of }A)$, a combination of the rows of $A$;
so the row space of $R$ is contained in that of $A$. Since $A=E^{-1}R$, each row of $A$ is likewise a
combination of the rows of $R$, giving the reverse inclusion. The zero rows of $R$ contribute
nothing, so the $r$ nonzero rows span. They are independent: the $i$th nonzero row has a $1$ in the
$i$th pivot column, where every other row of $R$ has a $0$, so in any combination equal to
$\vb{0}$ the $i$th coefficient is $0$.
\end{proof}

\begin{corollary}[Row rank equals column rank]\label{cor:u11-rankT}
For every matrix, $\dim\Col{A\T}=\dim\Col{A}$; that is, $\rank A\T=\rank A$.
\end{corollary}
\begin{proof}
Both dimensions equal the number $r$ of pivots of $A$, by Theorems~\ref{thm:u11-row}
and~\ref{thm:u10-col}. Applying Theorem~\ref{thm:u10-col} to $A\T$ shows that $\dim\Col{A\T}$ is
also the rank of $A\T$.
\end{proof}

This is a statement with content: there is no evident reason why the number of independent rows of
a $3\times 7$ matrix should equal its number of independent columns, yet it always does.

\subsection{The left nullspace and the matrix \texorpdfstring{\(E\)}{E}}

To find $\Nul{A\T}$ without transposing, we keep a record of the elimination. Attach the $m\times m$
identity to $A$ and apply to $\begin{bmatrix}A&I\end{bmatrix}$ the row operations that turn $A$
into $R$:
\[
\begin{bmatrix}A_{m\times n}&I_{m\times m}\end{bmatrix}\longrightarrow
\begin{bmatrix}R_{m\times n}&E_{m\times m}\end{bmatrix}.
\]
The row operations amount to multiplying on the left by some invertible matrix, and that matrix
turned $I$ into $E$, so it \emph{is} $E$: we get $EA=R$. When $A$ is square and invertible, $R=I$
and this is Gauss--Jordan: $E=A^{-1}$. In general $E$ is not unique when $r<m$ (one may add
multiples of the bottom rows to the others), but any such $E$ works below.

\begin{theorem}[Left nullspace]\label{thm:u11-left}
Let $EA=R$ with $E$ invertible and $R=\operatorname{rref}(A)$. The last $m-r$ rows of $E$ form a
basis of $\Nul{A\T}$, and $\dim\Nul{A\T}=m-r$.
\end{theorem}
\begin{proof}
$A\T$ is $n\times m$ with rank $r$ (Corollary~\ref{cor:u11-rankT}), so by
Theorem~\ref{thm:u10-null} its nullspace has dimension $m-r$. The last $m-r$ rows of $R=EA$ are
zero, so each of the last $m-r$ rows $\vb{y}\T$ of $E$ satisfies $\vb{y}\T A=\vb{0}\T$. These rows
are independent because $E$ is invertible (Theorem~\ref{thm:u10-test} applied to $E\T$). So they
are $m-r$ independent vectors in an $(m-r)$-dimensional space, and Proposition~\ref{prop:u10-right}(a)
makes them a basis.
\end{proof}

\begin{theorem}[Dimensions of the four subspaces]\label{thm:u11-dims}
For an $m\times n$ matrix of rank $r$:
\[
\dim\Col{A}=r,\qquad \dim\Nul{A}=n-r,\qquad \dim\Col{A\T}=r,\qquad \dim\Nul{A\T}=m-r.
\]
In $\R^n$: $r+(n-r)=n$. In $\R^m$: $r+(m-r)=m$.
\end{theorem}
\begin{proof}
Theorems~\ref{thm:u10-col}, \ref{thm:u10-null}, \ref{thm:u11-row} and~\ref{thm:u11-left}.
\end{proof}

The count $r+(n-r)=n$ is the old count of pivot variables plus free variables; the count in $\R^m$
is the same statement for $A\T$, which is just as good a matrix as $A$.

\begin{example}[All four subspaces at once]\label{ex:u11-four}
Take $A=\begin{bmatrix}1&2&3&1\\1&1&2&1\\1&2&3&1\end{bmatrix}$ from Example~\ref{ex:u10-CN},
so $m=3$, $n=4$, $r=2$. Reduce $\begin{bmatrix}A&I\end{bmatrix}$: subtract row~1 from rows~2 and~3,
multiply row~2 by $-1$, subtract $2\times$row~2 from row~1:
\[
\left[\begin{array}{cccc|ccc}1&2&3&1&1&0&0\\1&1&2&1&0&1&0\\1&2&3&1&0&0&1\end{array}\right]
\longrightarrow
\left[\begin{array}{cccc|ccc}1&0&1&1&-1&2&0\\0&1&1&0&1&-1&0\\0&0&0&0&-1&0&1\end{array}\right].
\]
Check $EA=R$ row by row: $-(1,2,3,1)+2(1,1,2,1)=(1,0,1,1)$, $(1,2,3,1)-(1,1,2,1)=(0,1,1,0)$,
$-(1,2,3,1)+(1,2,3,1)=\vb{0}$. Now read off:
\begin{itemize}
\item $\Col{A}$ (dimension $2$): pivot columns of $A$, $(1,1,1)$ and $(2,1,2)$.
\item $\Nul{A}$ (dimension $4-2=2$): special solutions $(-1,-1,1,0)$ and $(-1,0,0,1)$.
\item $\Col{A\T}$ (dimension $2$): nonzero rows of $R$, $(1,0,1,1)$ and $(0,1,1,0)$.
\item $\Nul{A\T}$ (dimension $3-2=1$): last row of $E$, $\vb{y}=(-1,0,1)$; indeed
$-(\text{row }1)+(\text{row }3)$ is the zero row.
\end{itemize}
\end{example}

\begin{trap}
Elimination keeps the row space and the nullspace but changes the column space and the left
nullspace. In Example~\ref{ex:u11-four}, $(1,1,1)$ is a column of $A$, so it lies in $\Col{A}$;
but every column of $R$ has third component $0$, so $(1,1,1)\notin\Col{R}$. Take the
\emph{rows of $R$} for the row space, but the \emph{columns of $A$} for the column space. Likewise
$\Nul{R\T}$ is spanned by $(0,0,1)$, not by the correct $(-1,0,1)$.
\end{trap}

\subsection{What a matrix keeps and what it crushes}

\begin{example}[The four subspaces of a rank-one matrix]\label{ex:u11-crush}
Return to $A=\begin{bmatrix}1&2&2\\2&4&4\end{bmatrix}$ from Example~\ref{ex:u08-crush}, with
$m=2$, $n=3$ and $r=1$.
\begin{itemize}
\item \emph{Inputs, $\R^3$.} The row space is the line through $(1,2,2)$ (dimension $r=1$). The
nullspace is the plane spanned by $\vb{s}_1=(-2,1,0)$ and $\vb{s}_2=(-2,0,1)$ (dimension
$n-r=2$). Both special solutions are perpendicular to $(1,2,2)$. Count: $1+2=3$.
\item \emph{Outputs, $\R^2$.} Every column is a multiple of $(1,2)$, so the column space is the
line through $(1,2)$ (dimension $1$). Row~2 minus $2\times$row~1 is zero, so $\vb{y}=(-2,1)$, or
equally $(2,-1)$, spans the left nullspace (dimension $m-r=1$). Count: $1+1=2$.
\end{itemize}
Now take the input $\vb{x}=(-1,3,2)$. It splits as a row-space part plus a nullspace part,
\[
\vb{x}=\begin{bmatrix}-1\\3\\2\end{bmatrix}=\begin{bmatrix}1\\2\\2\end{bmatrix}+\begin{bmatrix}-2\\1\\0\end{bmatrix},
\qquad
A\vb{x}=\begin{bmatrix}-1+6+4\\-2+12+8\end{bmatrix}=\begin{bmatrix}9\\18\end{bmatrix}
=A\begin{bmatrix}1\\2\\2\end{bmatrix}.
\]
The nullspace part was crushed to $\vb{0}$; only the row-space part reached the output, landing at
$(9,18)$ on the column-space line. The information that $A$ destroys is not scattered at random:
it is exactly the nullspace component, and that subspace is perpendicular to the rows.
(Section~\ref{sec:u16} proves that the nullspace is always perpendicular to the row space, and the
left nullspace to the column space; here $(1,2)\cdot(2,-1)=0$.) Figure~\ref{fig:u11-crush}
draws both sides.
\end{example}

\begin{figure}[ht]\centering
\begin{tikzpicture}[>=Stealth,font=\small]
  % input space R^3 (schematic): nullspace plane as a floor, row space as the vertical line
  \begin{scope}
    \fill[gray!12] (-2.3,-0.9) -- (1.5,-0.9) -- (2.5,0.55) -- (-1.3,0.55) -- cycle;
    \draw[gray!60!black] (-2.3,-0.9) -- (1.5,-0.9) -- (2.5,0.55) -- (-1.3,0.55) -- cycle;
    \node[gray!40!black,anchor=north] at (-0.4,-0.95) {nullspace plane, dim 2};
    \coordinate (O) at (0,-0.2);
    \draw[blue!60!black,dashed] (0,-0.9) -- (0,2.75);
    \node[blue!60!black,anchor=south] at (0,2.75) {row space, dim 1};
    \draw[->,thick,gray!50!black] (O) -- (1.35,0.05) node[below right=-1pt] {$\vb{s}_1$};
    \draw[->,thick,gray!50!black] (O) -- (-1.25,0.3) node[above left=-1pt] {$\vb{s}_2$};
    \draw[->,very thick,blue!70!black] (O) -- (0,2.0) node[left] {$(1,2,2)$};
    \draw[dashed] (0,2.0) -- (1.35,2.25) -- (1.35,0.05);
    \draw[->,very thick,orange!80!black] (O) -- (1.35,2.25) node[right] {$\vb{x}=(-1,3,2)$};
    \fill (O) circle (1.5pt);
    \node at (0.2,-1.75) {inputs: $\R^3$, \;$1+2=3$};
  \end{scope}
  % map
  \draw[->,thick] (3.3,1.0) -- node[above] {$A$} (4.7,1.0);
  % output space R^2, true proportions (scale 0.13)
  \begin{scope}[xshift=6.6cm,yshift=-0.3cm,scale=0.13]
    \draw[->,gray] (-14,0) -- (16,0);
    \draw[->,gray] (0,-9) -- (0,23);
    \draw[orange!80!black,thick] (-4.5,-9) -- (11,22) node[right] {column space, dim 1};
    \draw[gray!50!black,thick] (-15,7.5) -- (15,-7.5) node[below=1pt,pos=0.92] {left nullspace};
    \draw (1.8,3.6) -- (3.6,2.7) -- (1.8,-0.9);
    \fill[orange!80!black] (9,18) circle (9pt) node[right=3pt] {$A\vb{x}=(9,18)$};
    \fill (0,0) circle (7pt);
    \node[anchor=north east] at (-3,-1) {$A\vb{s}_1=\vb{0}$};
    \node at (2,-15) {outputs: $\R^2$, \;$1+1=2$};
  \end{scope}
\end{tikzpicture}
\caption{The matrix $A=\begin{bmatrix}1&2&2\\2&4&4\end{bmatrix}$. Left (schematic): the input
$\vb{x}=(-1,3,2)$ is the row-space vector $(1,2,2)$ plus the nullspace vector $\vb{s}_1$.
Right (to scale): $A$ crushes $\vb{s}_1$ to $\vb{0}$ and sends the row-space part to
$(9,18)$ on the column-space line, which is perpendicular to the left nullspace line through
$(2,-1)$.}
\label{fig:u11-crush}
\end{figure}

\subsection{A new vector space: matrices}

Nothing in the definition of a vector space requires the vectors to be columns of numbers. All we
use is that vectors can be added and multiplied by scalars, following the usual rules.

\begin{example}[The space of \texorpdfstring{\(3\times3\)}{3 by 3} matrices]\label{ex:u11-M}
Let $M$ be the set of all $3\times3$ real matrices. Matrices can be added and multiplied by
numbers, there is a zero matrix, and all eight vector-space rules hold; that we can also
\emph{multiply} two matrices is irrelevant here. Three subspaces of $M$: all upper triangular
matrices $U$; all symmetric matrices $S$; and their intersection $S\cap U$, which is the set $D$ of
diagonal matrices (symmetric with zeros below the diagonal forces zeros above it). Each is closed
under addition and scalar multiplication. A basis of $D$ is
\[
\begin{bmatrix}1&0&0\\0&0&0\\0&0&0\end{bmatrix},\quad
\begin{bmatrix}0&0&0\\0&3&0\\0&0&0\end{bmatrix},\quad
\begin{bmatrix}0&0&0\\0&0&0\\0&0&7\end{bmatrix}:
\]
they are independent (a combination with weights $a,b,c$ is $\operatorname{diag}(a,3b,7c)$, which
is zero only when $a=b=c=0$), and every diagonal matrix $\operatorname{diag}(d_1,d_2,d_3)$ is the
combination with weights $d_1,\,d_2/3,\,d_3/7$. So $\dim D=3$. Section~\ref{sec:u12} finds the
dimensions of $M$, $S$ and $U$.
\end{example}

\subsection{Exercises}

\begin{exercise}\label{exr:u11-ineq}
$A$ is $m\times n$ of rank $r$, and some right-hand sides $\vb{b}$ give no solution of
$A\vb{x}=\vb{b}$. (a) Which inequalities ($<$ or $\le$) must hold between $m$, $n$ and $r$?
(b) How do you know that $A\T\vb{y}=\vb{0}$ has solutions other than $\vb{y}=\vb{0}$?
\end{exercise}
\begin{solution}
(a) Always $r\le m$ and $r\le n$. Since some $\vb{b}$ is unreachable, $\Col{A}\ne\R^m$, so
$r<m$ (Theorem~\ref{thm:u09-fullrank}(b)). Nothing more can be said about $r$ versus $n$.
(b) $\dim\Nul{A\T}=m-r>0$, so the left nullspace contains nonzero vectors.
\end{solution}

\begin{exercise}\label{exr:u11-ATy}
For which $\vb{d}$ is $A\T\vb{y}=\vb{d}$ solvable? When is the solution unique?
\end{exercise}
\begin{solution}
$A\T\vb{y}$ is a combination of the columns of $A\T$, i.e.\ of the rows of $A$, so the system is
solvable exactly when $\vb{d}$ is in the row space $\Col{A\T}$. Two solutions differ by a vector in
$\Nul{A\T}$, so the solution is unique when the left nullspace is $\{\vb{0}\}$, i.e.\ when $r=m$.
\end{solution}

\begin{exercise}\label{exr:u11-LU}
Let $B=LU$ with $L=\begin{bmatrix}1&0&0\\2&1&0\\-1&0&1\end{bmatrix}$ and
$U=\begin{bmatrix}5&0&3\\0&1&1\\0&0&0\end{bmatrix}$. Find a basis and the dimension of each of the
four subspaces of $B$, without multiplying out $B$.
\end{exercise}
\begin{solution}
$U$ has two pivots, so $r=2$ and all four dimensions are $2$, $3-2=1$, $2$, $3-2=1$.
\emph{Column space:} since $L$ is invertible, the columns of $B=LU$ satisfy the same relations as
the columns of $U$, so columns~1 and~2 of $B$ are pivot columns. They are $L$ times $(5,0,0)$ and
$(0,1,0)$, i.e.\ $5(1,2,-1)$ and $(0,1,0)$; a basis is $(1,2,-1),(0,1,0)$, the first two columns
of $L$. \emph{Nullspace:} $\Nul{B}=\Nul{U}$; the free variable $x_3=1$ gives $x_2=-1$ and
$5x_1+3=0$, so the basis is $(-\tfrac35,-1,1)$. \emph{Row space:} $L^{-1}B=U$ is a sequence of row
operations, so the row space is spanned by the pivot rows of $U$: $(5,0,3)$ and $(0,1,1)$.
\emph{Left nullspace:} $L^{-1}=\begin{bmatrix}1&0&0\\-2&1&0\\1&0&1\end{bmatrix}$ plays the role of
$E$, and its last row $(1,0,1)$ is the basis: row~1 of $B$ is $(5,0,3)$, row~3 is $-(5,0,3)$, and
they cancel.
\end{solution}

\begin{exercise}\label{exr:u11-construct}
Find a $3\times3$ matrix whose column space contains $(1,1,0)$ and $(1,0,1)$ and whose nullspace
contains $(1,1,1)$. What must its rank be?
\end{exercise}
\begin{solution}
The column space contains two independent vectors, so $r\ge2$; the nullspace is at least a line,
so $3-r\ge1$. Hence $r=2$. The condition $A(1,1,1)=\vb{0}$ says the three columns add to zero. Take
the first two columns to be the given vectors and the third their negative sum:
$A=\begin{bmatrix}1&1&-2\\1&0&-1\\0&1&-1\end{bmatrix}$. Its column space is the plane spanned by
$(1,1,0)$ and $(1,0,1)$, and $A(1,1,1)=(0,0,0)$.
\end{solution}
